\chapter*{\textit{Conclusion}}\addcontentsline{toc}{chapter}{\textit{Conclusion}}
Au terme de cet exposé, nous avons parcouru les concepts fondamentaux qui structurent la gestion des données dans un système informatique.

Nous avons établi qu'un \textbf{système de fichiers} est bien plus qu'un simple espace de stockage~: c'est une structure de données logique qui offre à l'utilisateur une vue abstraite et cohérente de ses informations, organisées sous forme d'arborescence de fichiers et de répertoires.

L'étude de l'\textbf{organisation hiérarchique} a mis en évidence une différence fondamentale d'approche entre Windows, qui expose chaque partition comme un lecteur indépendant, et Unix/Linux, qui intègre l'ensemble des systèmes de fichiers dans une arborescence unique grâce au mécanisme de \textbf{montage}. Ce dernier, piloté par les commandes \texttt{mount} et \texttt{umount} et configuré via \texttt{/etc/fstab}, constitue un mécanisme central de l'administration système sous Unix.

Enfin, la diversité des \textbf{types de systèmes de fichiers} locaux (ext2, FAT, NTFS), réseau (NFS, SMB), optiques (ISO~9660, UDF) et journalisés (ext3, ext4, XFS) reflète la variété des contextes d'utilisation et des exigences en matière de performance, de compatibilité et de fiabilité. La \textbf{journalisation} en particulier représente une avancée majeure pour la cohérence des données en cas de panne.

La maîtrise des systèmes de fichiers est ainsi une compétence essentielle pour tout professionnel de l'informatique, à la croisée de l'administration système, de la sécurité des données et de l'architecture des systèmes de stockage.